\documentclass[../../../include/open-logic-section]{subfiles}

\begin{document}
\olfileid[tr]{sth}{card-arithmetic}{simp}

\olsection{Toplama ve Çarpmanın Sadeleştirilmesi}

Sonlu ötesi kardinal toplaması ve çarpmasının \emph{son derece} kolay
olduğu ortaya çıkar. Bu, kardinal sayıların (belirli) sıra sayıları
olmasından, dolayısıyla iyi sıralanmasından ve belirli bir biçimde
işlenebilmesinden kaynaklanır. Ne var ki bunu göstermek \emph{o kadar da}
kolay değildir. Başlangıçta biraz incelikli bir tanıma ihtiyacımız var:

\begin{defn}
Sıra sayısı çiftleri üzerinde bir \emph{kanonik sıralama} $\canonord$
tanımlarız: $\tuple{\alpha_1, \alpha_2} \canonord
\tuple{\beta_1, \beta_2}$ ancak ve ancak aşağıdakilerden biri geçerliyse:
\begin{enumerate}
	\item $\max(\alpha_1, \alpha_2) < \max(\beta_1, \beta_2)$; ya da
	\item $\max(\alpha_1, \alpha_2) = \max(\beta_1, \beta_2)$ ve
	$\alpha_1 < \beta_1$; ya da
	\item $\max(\alpha_1, \alpha_2) = \max(\beta_1, \beta_2)$ ve
	$\alpha_1 = \beta_1$ ve $\alpha_2 < \beta_2$.
\end{enumerate}
\end{defn}

\begin{lem}
Her $\alpha$ sıra sayısı için $\tuple{\alpha \times \alpha, \canonord}$
bir iyi sıralamadır.
\end{lem}

\begin{proof}
$\canonord$ bağıntısının $\alpha\times\alpha$ üzerinde bağlantılı olduğu
açıktır. Ne $\tuple{\alpha_1,\alpha_2}$ ne de
$\tuple{\beta_1,\beta_2}$ ötekinden $\canonord$ bakımından küçük olsun.
O zaman $\max(\alpha_1,\alpha_2)=\max(\beta_1,\beta_2)$,
$\alpha_1=\beta_1$ ve $\alpha_2=\beta_2$ olur; dolayısıyla
$\tuple{\alpha_1,\alpha_2}=\tuple{\beta_1,\beta_2}$.

İyi sıralamayı göstermek için boş olmayan bir
$X\subseteq\alpha\times\alpha$ alın. $\alpha$ bir sıra sayısı olduğundan,
bir $\delta$, $\Setabs{\max(\gamma_1, \gamma_2)}{\tuple{\gamma_1,
\gamma_2} \in X}$ kümesinin en küçük elemanıdır. Şimdi
$\Setabs{\tuple{\gamma_1,\gamma_2} \in X}{\max(\gamma_1, \gamma_2) =
\delta}$ kümesindeki ilk bileşeni en küçük olanlar dışındaki bütün
çiftleri atın; kalanların arasında ikinci bileşeni en küçük olan çift,
$X$'in $\canonord$-en küçük elemanıdır.
\end{proof}
\noindent
Şimdi küçük ve basit bir gözlem:

\begin{prop}\ollabel{simplecardproduct}
$\cardeq{\alpha}{\beta}$ ise
$\cardeq{\alpha \times \alpha}{\beta \times \beta}$.
\end{prop}

\begin{proof}
$f\colon\alpha\to\beta$ fonksiyonunun indüklediği
$\tuple{\gamma_1,\gamma_2}\mapsto
\tuple{f(\gamma_1),f(\gamma_2)}$ eşlemesini alın.
\end{proof}
\noindent
Şimdi bütün bunları temel bir lemayı ispatlamak için kullanacağız:
\begin{lem}\ollabel{alphatimesalpha}
Her sonsuz $\alpha$ sıra sayısı için
$\cardeq{\alpha}{\alpha\times\alpha}$.
\end{lem}

\begin{proof}
Olmayana ergiyle bunun yanlış olduğu en küçük sonsuz sıra sayısı $\alpha$
olsun. \olref[sfr][siz][zigzag]{natsquaredenumerable},
$\cardeq{\omega}{\omega\times\omega}$ olduğunu gösterir; dolayısıyla
$\omega\in\alpha$. Üstelik $\alpha$ bir kardinal sayıdır: olmayana
ergiyle olmadığını varsayın. O zaman $\card{\alpha}\in\alpha$ olur;
varsayım gereği
$\cardeq{\card{\alpha}}{\card{\alpha}\times\card{\alpha}}$ ve tanım
gereği $\cardeq{\card{\alpha}}{\alpha}$ olur; dolayısıyla
\olref{simplecardproduct} gereği
$\cardeq{\alpha}{\alpha\times\alpha}$ elde edilir.

Şimdi her $\tuple{\gamma_1,\gamma_2}\in\alpha\times\alpha$ için şu
kesiti ele alın:
\begin{align*}
	\text{Seg}(\gamma_1, \gamma_2) &= \Setabs{\tuple{\delta_1, \delta_2} \in \alpha \times \alpha}{\tuple{\delta_1, \delta_2} \canonord \tuple{\gamma_1, \gamma_2}}
\end{align*}
$\gamma=\max(\gamma_1,\gamma_2)$ olsun. O zaman
$\tuple{\gamma_1,\gamma_2}\canonord
\tuple{\gamma+1,\gamma+1}$ olduğuna dikkat edin. Dolayısıyla $\gamma$
sonsuz olduğunda:
\begin{align*}
	\text{Seg}(\gamma_1, \gamma_2) &
	\precsim ((\gamma \ordplus 1)\ordtimes (\gamma \ordplus 1))\\
	&\approx (\gamma \ordtimes \gamma)
	\text{, \olref[ord-arithmetic][using-addition]{ordinfinitycharacter} ve
	\olref{simplecardproduct} gereği}\\
	&\approx \gamma \text{, tümevarım varsayımı gereği}\\
	& \prec \alpha\text{, çünkü $\alpha$ bir kardinal sayıdır}
\end{align*}
olur. Böylece $\ordtype{\alpha\times\alpha,\canonord}\leq\alpha$ ve
dolayısıyla $\cardle{\alpha\times\alpha}{\alpha}$ olur. Elbette
$\cardle{\alpha}{\alpha\times\alpha}$ olduğundan sonuç
Schröder--Bernstein ile çıkar.
\end{proof}

Nihayet sadeleştirme sonucumuza geliyoruz:

\begin{thm}\ollabel{cardplustimesmax}
$\cardfont{a},\cardfont{b}$ sonsuz kardinal sayılarsa:
\[
	\cardfont{a}
\cardtimes \cardfont{b} = \cardfont{a} \cardplus \cardfont{b} =
\text{max}(\cardfont{a}, \cardfont{b}).
\]
\end{thm}

\begin{proof}
Genelliği kaybetmeden $\cardfont{a}=\max(\cardfont{a},\cardfont{b})$
varsayın. O zaman \olref{alphatimesalpha} uygulanarak
$\cardfont{a}\cardtimes\cardfont{a}=\cardfont{a}\leq\cardfont{a}
\cardplus\cardfont{b}\leq\cardfont{a}\cardplus\cardfont{a}\leq
\cardfont{a}\cardtimes\cardfont{a}$ elde edilir.
\end{proof}\noindent
Benzer biçimde $\cardfont{a}$ sonsuzsa, her biri
$\leq\cardfont{a}$ büyüklüğündeki $\cardfont{a}$ tane kümenin birleşimi
$\leq\cardfont{a}$ büyüklüğündedir:

\begin{prop}\ollabel{kappaunionkappasize}
$\cardfont{a}$ sonsuz bir kardinal sayı olsun. Her
$\beta\in\cardfont{a}$ sıra sayısı için $X_\beta$,
$\card{X_\beta}\leq\cardfont{a}$ olan bir küme olsun. O zaman
$\card{\bigcup_{\beta \in \cardfont{a}} X_\beta}\leq\cardfont{a}$.
\end{prop}

\begin{proof}
Her $\beta\in\cardfont{a}$ için bir $f_\beta\colon
X_\beta\to\cardfont{a}$ !!a{injection} sabitleyin.\footnote{Bunlar nasıl
``sabitlenir''? Bkz. \olref[sth][choice][countablechoice]{sec}.}
$v\in X_\beta$ ve hiçbir $\gamma\in\beta$ için $v\notin X_\gamma$
olmamak üzere $g(v)=\tuple{\beta,f_\beta(v)}$ ile bir
$g\colon\bigcup_{\beta \in \cardfont{a}}X_\beta\to
\cardfont{a}\times\cardfont{a}$ !!a{injection} tanımlayın. Şimdi
\olref{cardplustimesmax} gereği
$\bigcup_{\beta \in \cardfont{a}} X_\beta \preceq
\cardfont{a}\times\cardfont{a}\approx\cardfont{a}$.
\end{proof}

\end{document}
